% Conclusion and bibliography of the FeRRy paper (draft, 8 Oct 2026). Appended by
% DeveloperDocs/paper/assemble.py after evaluation.tex.
%%%%% BEGIN SECTION %%%%%
\section{Conclusion and Future Work}
\label{sec:conclusion}

This paper presented \sysname{}, a framework that sustains federated learning
across disconnected regions by planning UAV data-mule missions in which the
radio's reach is a decision. At the dock, \sysname{} chooses a contact band
class and a route together under one feasibility test that respects per-device
deadlines and the mission budget. In flight, the same test gates a per-stop
decision and a re-plan, and staleness-weighted merges on the mule and at the
server close rounds with the updates that arrive. In a multi-process prototype
that trains a real intrusion-detection model, \sysname{} reached the target
accuracy 27--31\% sooner than seven whole schedulers at six devices, because its
plan chose a long-reach narrow band that removed most of the transit between
stops. It kept updates fresher than age-of-information schedulers under tight
budgets, beat the FedEx route at 24 devices under the stress budget, and planned
in 1.5\,s at 96 devices. Several elements of the design showed no measurable
effect at our settings. Among them is a learned per-stop score: FQ performed like
the fixed rule, and its exploratory advantage over a learned scheduler without a
plan came from the plan it flies inside.

Future work follows from the limitations we measured. \sysname{}'s plan prices
the mean channel, so it overran tight budgets in 30\% of missions; pricing the
channel's variation, or reserving slack for it, is a first step. The plan also
ignores training time, so when devices need time to train it reaches devices
whose updates are not ready; adding a compute term to the plan is a second.
Beyond these, mission time, the radio channel and energy are simulated in this
evaluation. Validating the band trade over the air on a UAV testbed, with
measured channels and battery limits, and coordinating mules that share
spectrum and a dock remain open.

% IPDPS 2027 asks authors to disclose AI-generated text in the Acknowledgements.
% DRAFT wording: the authors must confirm or rewrite it. Funding acknowledgements
% stay out until the camera-ready (double-anonymous review).
\section*{Acknowledgment}
Parts of this paper's text, and of the code for its experiments and analysis,
were drafted with the assistance of an AI model (Claude, Anthropic). The authors
reviewed, edited, and verified all of it, and take responsibility for the
content.

\bibliographystyle{ieeetr}
\bibliography{biblio}
%%%%% END SECTION %%%%%
